\chapter{Acids and Bases: Beyond Sour and Slippery}

\begin{kaichang}
\textbf{Translator safety note:} Do not use taste, prolonged skin contact, or direct smelling to identify chemicals. The concentrated-acid and acid/ammonia examples require a trained teacher's risk assessment, not home imitation. Household does not mean harmless. See \href{https://www.acs.org/content/dam/acsorg/about/governance/committees/chemicalsafety/safetypractices/safety-in-the-elementary-school-science-classroom.pdf}{ACS classroom safety guidance}.

Cut a slice of lemon and touch it with your tongue: a sharp sourness immediately makes you squint. Now try something else: when washing your hands, deliberately rub the soap a little longer and notice the slippery feeling between your fingers that seems impossible to rinse away.

Guess two things. First, is sourness the lemon's ``flavor,'' or is your tongue reporting a microscopic event? Second, is soapy water slippery because soap itself is slippery, or because it is doing something to your skin?

Here is an even less intuitive question: put dry ice, solid carbon dioxide, into water, and the water becomes acidic; add aqueous ammonia, and it becomes ``basic.'' Yet dry ice and ammonia molecules contain neither hydrogen like hydrochloric acid nor hydroxide like caustic soda. How do they manage it?

By this chapter's end, you will see an unending proton relay behind the words ``acid'' and ``base.'' First, write down your guesses about lemon and soap.
\end{kaichang}

\section{Acids on the Tongue, Bases on the Skin}

Human knowledge of acids and bases began with the body. Vinegar, lemons, sour plums, yogurt: acidic things taste sour. Ancient people even named vinegar for its sourness; the English word acid comes from Latin for ``sour.'' Bases took a more roundabout path: people noticed substances that removed grease and felt slippery---water steeped with plant ash, soapy water, and baking-soda solution. Arabic speakers called the active material in plant ash al-qaly, the root of the English word alkali.

With only tongue and hands, people assembled their first ``files'' on acids and bases:

\begin{itemize}[itemsep=2pt]
\item Most acidic solutions taste sour. \textbf{This is everyday experience, never permission to taste anything in a laboratory}: many acids burn the mouth, and many have toxicity unrelated to sourness.
\item Acids change the color of certain plant juices: purple cabbage and morning-glory juice turn red in acid.
\item Acids bubble on limestone, shells, and baking soda, releasing carbon dioxide.
\item Basic solutions feel slippery, remove grease, and turn purple cabbage juice green or even yellow.
\item Mixed acids and bases cancel each other's ``temperaments'': pour vinegar into baking-soda solution, and after the fizzing subsides it is neither sour nor slippery.
\end{itemize}

For centuries, people repeatedly used these effects: cleaning metals with acid, making soap with plant ash, and neutralizing acid to ``put out a fire'' in the stomach. But files are not understanding. Why do all acids---hydrochloric, acetic, citric, carbonic---produce the same sour sensation despite their completely different molecules?

The simplest explanation for different substances producing the same sensation is that they release \textbf{the same kind of particle} in water.

\section{The First Definition: What They Release in Water}

In the 1880s, a young Swedish chemistry doctoral student focused on this possibility. Remember the conclusion first; the story comes shortly. His name was Arrhenius, and he proposed:

\begin{qiaoliang}
\textbf{The Arrhenius definition}: an acid releases hydrogen ions, \chem{H^+}, in water; a base releases hydroxide ions, \chem{OH^-}, in water. Neutralization combines \chem{H^+} and \chem{OH^-} to make water.
\end{qiaoliang}

This immediately explained why acids all taste sour: hydrochloric, acetic, and citric acids all release \chem{H^+} in water, and your taste buds essentially detect its concentration. It also explained neutralization: \chem{H^+} and \chem{OH^-} meet and embrace to form water molecules, making both ``disappear.''

But we must immediately correct the particle picture. Almost no naked \chem{H^+} exists in aqueous solution. When a hydrogen atom loses its sole electron, what remains is a \textbf{bare proton}, a tiny charged particle with a formidable electric field. Chapter 23 explained that water molecules are polar, with negatively charged oxygen ends. A bare proton is therefore immediately ``embraced'' by nearby oxygen ends, forming hydronium, \chem{H_3O^+}, which also carries other water molecules along. The more honest notation is:

\[\chem{HCl} + \chem{H_2O} \rightarrow \chem{H_3O^+} + \chem{Cl^-}\]

\[\chem{NaOH} \rightarrow \chem{Na^+} + \chem{OH^-}\]

For convenience, people usually write \chem{H^+}; you will often see this shorthand too. Just remember that the ``protons'' roaming acidic solutions ride on water molecules.

Arrhenius's definition is neat, but cannot answer the opening question: ammonia, \chem{NH_3}, dissolves to give a basic solution that neutralizes acids and changes indicators. Yet an ammonia molecule has no oxygen atom at all: where could it get \chem{OH^-}? Furthermore, bring the mouths of hydrogen chloride and ammonia gas bottles together, and they react directly to form white smoke---tiny ammonium chloride particles, \chem{NH_4Cl}. No water is involved, so Arrhenius's definition would not even classify this as an acid--base reaction, although it clearly looks like acid and base canceling each other.

A useful definition that cannot explain new facts signals that a theory needs upgrading.

\begin{zhengju}
How did Arrhenius think of ``dissociation''? Not by tasting, but through two sets of puzzling experimental data.

The first concerned conductivity. Pure water conducts very poorly, but adding salt, acid, or base immediately improves it. The explanation then was that electricity split molecules into charged fragments. Arrhenius noticed something awkward: conductivity's dependence on concentration did not look like fragments produced by electricity, but like charged particles already present whether electricity flowed or not.

The second set was more compelling. Chapter 25 explained that freezing-point depression and osmotic pressure count only particle numbers. Experiments found that dissolving \ud{1}{mol} of salt lowered water's freezing point about twice as much as \ud{1}{mol} of sugar; calcium chloride, \chem{CaCl_2}, produced nearly three times the effect. If dissolved salts remained intact molecular units, why were two or three sets of particles being counted? Leading chemists, including Mendeleev, offered alternatives: perhaps solute and water formed loose ``hydrates.'' The debate remained unresolved.

In his 1884 doctoral thesis, Arrhenius proposed a bold unified answer: salts, acids, and bases \textbf{split into charged ions by themselves} in water, without electricity. Salt becomes two kinds of particles, \chem{Na^+} and \chem{Cl^-}, so freezing-point depression counts two; hydrochloric acid becomes \chem{H^+} and \chem{Cl^-}, explaining conductivity and sourness. The examining professors thought this unknown young man's idea fanciful and awarded his thesis a low, barely passing grade.

Arrhenius refused to accept defeat. He sent the thesis to several European chemists and gained support from van 't Hoff, Ostwald, and others. Crucial reinforcement came from physics: Thomson discovered the electron in 1895, making ``charged atoms'' less absurd, since atoms could lose or gain electrons. In 1903, that barely passing thesis won Arrhenius the Nobel Prize in Chemistry. What matters is not the dramatic vindication but the method. He had no microscope capable of seeing ions; he relied on \textbf{several independent phenomena pointing to one explanation}. Conductivity, freezing points, sourness, and neutralization all required ions already present in solution. When several lines of evidence converge on the same unseen particle, we dare to give it an identity card.
\end{zhengju}

\section{A Broader Definition: Protons Change Hands}

In 1923, the Dane Brønsted and the Englishman Lowry independently expanded Arrhenius's definition. Their idea was to stop focusing on which ions are released and follow \textbf{the proton itself}.

\begin{qiaoliang}
\textbf{The Brønsted definition}: an acid \textbf{donates a proton} (\chem{H^+}); a base \textbf{accepts a proton}. An acid--base reaction transfers a proton from one particle to another, and need not occur in water.
\end{qiaoliang}

Seen this way, aqueous ammonia makes sense. Ammonia does not release its own \chem{OH^-}; it \textbf{takes} a proton from a water molecule:

\[\chem{NH_3} + \chem{H_2O} \rightleftharpoons \chem{NH_4^+} + \chem{OH^-}\]

Ammonia is the base, accepting a proton; water acts as the acid, donating one. The water molecule, left with only half its body after the theft, becomes \chem{OH^-}. Notice the reversible arrow: this is Chapter 31's dynamic equilibrium. After most ammonia molecules take a proton, the resulting ammonium, \chem{NH_4^+}, gives it back. At any instant, only a small fraction is present as \chem{NH_4^+} and \chem{OH^-}. That is why aqueous ammonia is a ``mild'' base.

Brønsted's definition also reveals water's dual personality. With hydrochloric acid, water accepts a proton to become \chem{H_3O^+}, acting as a base; with ammonia, it donates a proton to become \chem{OH^-}, acting as an acid. Water is both, depending on its partner. Water molecules even exchange protons with one another: one hands a proton to a neighbor, becoming \chem{OH^-}, while the neighbor becomes \chem{H_3O^+}. This is water's autoionization. At every moment, a few parts per billion of pure water's molecules are in this ``split'' state. Is that negligible? Quite the opposite: the whole discussion of pH later in this chapter rests on it.

A still more elegant consequence follows. After donating a proton, the remainder of an acid can take it back, so a base stands behind every acid; after accepting a proton, a base can donate it, so an acid stands behind every base. Such partners are \textbf{conjugate acid--base pairs}: acetate is acetic acid's conjugate base, and ammonia is ammonium's. The stronger an acid, the harder it is to reclaim its donated proton and the weaker its conjugate base---like an echo growing ever fainter.

The answer for dry ice now emerges too: \chem{CO_2} itself donates no proton, but first reacts with water to form carbonic acid, \chem{H_2CO_3}, which then donates one. This route---dissolved \chem{CO_2} making water acidic---sounds like classroom trivia, yet determines the fate of oceans and blood. We will return to it later.

For those looking further ahead, there is a third and broadest definition: the \textbf{Lewis definition}. An acid accepts an electron pair, and a base donates one. It encompasses Brønsted's definition---ammonia donating its lone pair to a proton is a Lewis base---and also explains reactions with no protons involved, such as water or ammonia molecules surrounding a metal ion to form a complex ion. That belongs to Chapter 35's coordination chemistry. For now, this is a signpost: three successive definitions, each broader than its predecessor, enlarging the map rather than overturning what came before. In solutions, Brønsted's ``protons changing hands'' remains the most convenient tool.

\section{Strong and Concentrated Are Different}

Next come two terms people confuse throughout their lives. First, fix this boundary firmly: \textbf{strength describes character; concentration describes amount. These are entirely independent concepts.}

An acid's strength describes \textbf{how completely} it donates protons in water. Hydrochloric acid is strong: almost every \chem{HCl} molecule entering water splits into \chem{H_3O^+} and \chem{Cl^-}, leaving none behind. Acetic acid is weak: once \chem{CH_3COOH} enters water, only about one percent of its molecules donate protons, while the other ninety-nine percent remain intact. Its reaction reaches equilibrium partway along:

\[\chem{CH_3COOH} + \chem{H_2O} \rightleftharpoons \chem{CH_3COO^-} + \chem{H_3O^+}\]

Concentration means something else entirely: \textbf{how much} acid you put into a liter of water. Dissolve \ud{10}{mol} of acetic acid in a liter of water, and it is concentrated acetic acid; dissolve \ud{0.0001}{mol} of hydrochloric acid in a liter, and it is dilute hydrochloric acid. All four combinations therefore exist, and all occur around us:

\begin{table}[htbp]
\centering
\begin{tabular}{lll}
\toprule
 & \shortstack[l]{Concentrated\\(much solute)} & \shortstack[l]{Dilute\\(little solute)} \\
\midrule
\shortstack[l]{Strong\\(complete ionization)} & \shortstack[l]{Concentrated sulfuric\\and hydrochloric acids\\(laboratory reagent bottles)} & \shortstack[l]{Stomach hydrochloric acid\\(about \ud{0.03}{mol\,L^{-1}})} \\
\shortstack[l]{Weak\\(incomplete ionization)} & \shortstack[l]{Glacial acetic acid\\(vinegar's concentrated\\ancestor)} & \shortstack[l]{White vinegar\\(about 5\% aqueous\\acetic acid)} \\
\bottomrule
\end{tabular}
\end{table}

Confusing these concepts can cause real danger. Hearing that acetic acid is ``weak,'' someone might assume glacial acetic acid is safe. Wrong: it is a concentrated weak acid and still burns skin. ``Weak'' means only that it is reluctant to donate protons; sheer quantity can overcome that reluctance. Conversely, extremely dilute hydrochloric acid is strong yet dilute enough to exist in your stomach. Weak does not mean gentle, and strong does not mean dangerous. The danger lies in \textbf{the total amount at a particular concentration}, the same principle as Chapter 1's ``the dose makes the poison.''

Remember one inviolable laboratory rule for later: when diluting concentrated sulfuric acid, always slowly add acid to water while stirring, never the reverse. Sulfuric acid releases much heat on meeting water, as in Chapter 27's energy accounting. Water poured into acid can boil at the surface and splatter hot acid into your face.

\begin{wuqu}
``Anything below pH 7 is corrosive and untouchable; alkaline things are gentle, slippery, and safe.'' Both halves are wrong. Against the first: gastric juice has pH about 1.5 and cola about 2.5, yet one is inside you and you drink the other. The second half is more dangerous: strong bases can be more insidiously corrosive than strong acids. Acid burns sting immediately, making you withdraw, whereas strong bases, such as caustic-soda drain cleaner, damage skin with \textbf{relatively little pain}. The slippery feeling actually comes from the base breaking surface skin oils down into soap---saponification. By the time you notice something wrong, damage has occurred. Always wear gloves with strongly alkaline products such as drain or oven cleaners, and never use slipperiness to ``identify'' a chemical. Judge acidity and basicity with indicators and pH meters, not skin or tongue.
\end{wuqu}

\section{Logarithmic Accounting for Protons}

Now let us quantify acidity. The concentration of \chem{H^+}---really \chem{H_3O^+}, remember---spans a staggering range: about \ud{10}{mol} per liter in concentrated hydrochloric acid, down to $10^{-15}$ mol per liter in concentrated caustic soda, sixteen orders of magnitude. To keep accounts across trillionfold ranges, logarithms are the natural tool: look not at the number itself, but at its power of ten.

\begin{qiaoliang}
\[\mathrm{pH} = -\lg\,[\chem{H^+}]\]
Here $[\chem{H^+}]$ is moles of hydrogen ions per liter of solution, strictly their ``activity,'' as the challenge explains. A pH difference of 1 means a tenfold hydrogen-ion concentration difference; a difference of 2 means a hundredfold difference. Higher concentration means lower pH because of the minus sign.
\end{qiaoliang}

Apply this ruler to the world around you: lemon juice is about pH 2.4, cola 2.5, white vinegar 2.9, milk 6.5, pure water at $25\,^{\circ}\mathrm{C}$ 7.0, blood 7.4, soapy water about 10, and drain cleaner can exceed 13.

An order-of-magnitude estimate shows the logarithmic scale's drama. Stomach acid at about pH 1.5 differs from blood at about 7.4 by 5.9 units. Hydrogen-ion concentrations differ by $10^{5.9}$, about eight hundred thousand: \textbf{a liter of stomach fluid contains nearly a million times as many hydrogen ions as a liter of blood}. Both fluids occupy the same body, separated only by the stomach wall. Why does that wall not digest itself? Mucus and rapidly renewed epithelial cells provide defenses; when these fail, a gastric ulcer results. In the other direction, seawater pH has fallen from about 8.2 before the Industrial Revolution to about 8.1 today, just 0.1 unit, but that means about $26\%$ more hydrogen ions ($10^{0.1}\approx1.26$). On a logarithmic scale, 0.1 is not trivial. We will return to the ocean's bill for that change.

\section{pH 7 Is Not a Universal Constant}

Almost every middle-school student has memorized ``pH 7 is neutral.'' Please revise that: \textbf{it holds only near $25\,^{\circ}\mathrm{C}$, not universally}.

Return to neutrality's true definition: not ``pH equals a certain number,'' but $[\chem{H^+}]=[\chem{OH^-}]$, equal hydrogen-ion and hydroxide-ion concentrations. Pure water always carries this balance: molecules exchange protons to produce pairs of \chem{H_3O^+} and \chem{OH^-} at equal concentrations. Thus \textbf{pure water is neutral at every temperature}, without exception.

But the number at which they are equal drifts with temperature. Chapter 31 taught that equilibrium shifts to counter an external change. Water's autoionization is endothermic---splitting a water molecule into two ions costs energy---so heating shifts equilibrium toward ionization, increasing both \chem{H^+} and \chem{OH^-}. Quantitatively, at $25\,^{\circ}\mathrm{C}$, pure water has $[\chem{H^+}]=10^{-7}$ mol/L, exactly pH 7.0. At $100\,^{\circ}\mathrm{C}$, $[\chem{H^+}]$ rises to about $10^{-6.1}$ mol/L, moving neutrality to about pH 6.1. At $0\,^{\circ}\mathrm{C}$, neutrality shifts to about pH 7.5.

Pure water at $100\,^{\circ}\mathrm{C}$ therefore has pH 6.1 yet remains neutral, because \chem{H^+} and \chem{OH^-} are still equally abundant. Conversely, a solution adjusted to exactly pH 7.0 at $100\,^{\circ}\mathrm{C}$ is \textbf{basic}: fewer hydrogen ions than the neutral value at that temperature means more hydroxide. Memorizing ``7 means neutral'' mistakes one thermometer mark for a physical law. The unchanging rule defines neutrality by the relative amounts of the two ions; pH reads only the \chem{H^+} side.

\begin{shiyan}
\textbf{Purple Cabbage: A Kitchen pH Meter}

Materials: a few purple cabbage leaves; hot water; five clear cups; white vinegar; baking-soda solution; soapy water; Sprite or another colorless carbonated drink; purified water. Ask an adult to help pour the hot water.

Procedure:
\begin{enumerate}[itemsep=1pt]
\item Tear up the cabbage leaves, steep them in hot water for ten minutes to obtain a deep-purple juice, then let it cool. Molecules called anthocyanins produce the purple color; their structures change with the number of protons they hold, changing their color too.
\item Put a little cabbage juice in each cup and add small amounts of vinegar, baking-soda solution, soapy water, Sprite, and purified water, respectively.
\item Record the colors and arrange them in red--purple--blue--green--yellow order.
\end{enumerate}

What to observe: the vinegar cup should be reddish, Sprite pinkish, purified water purple, baking soda blue-green, and soap green. One indicator gives five colors, showing that these liquids' \chem{H^+} concentrations are ordered by powers of ten. You have just made a logarithmic measurement by eye.

Safety: all materials are foods or household products and are safe, but \textbf{discard every experimental liquid without drinking it}: laboratory vessels are not tableware. Avoid hot-water burns. For an extension, try morning-glory juice, purple-sweet-potato water, or black tea, which changes color with lemon.
\end{shiyan}

\section{Three Tools: Neutralization, Buffers, Titration}

Three tools bring acid--base knowledge into practical use.

\textbf{Neutralization} is the most direct: \chem{H_3O^+} meets \chem{OH^-}, producing water and releasing heat. Mix equal amounts of dilute hydrochloric acid and dilute caustic soda solution, and the temperature rises noticeably, a classic entry in Chapter 27's energy ledger:

\[\chem{H_3O^+} + \chem{OH^-} \rightarrow \chem{2H_2O}\]

Antacids for excess stomach acid apply neutralization: sodium bicarbonate, calcium carbonate, and aluminum hydroxide, \chem{Al(OH)_3}, each use basicity to bring excess \chem{H^+} under control. Remember dose: a little for occasional heartburn is fine, but prolonged reliance can mask the real cause and disrupt the stomach's acid--base rhythm. Leave such questions to a doctor, rather than acting as your own pharmacist.

\textbf{Buffers} are subtler. Neutralization says ``eliminate whatever arrives''; buffering says ``however much arrives, try to remain unruffled.'' Consider blood's carbonic acid--bicarbonate buffer, a conjugate acid--base pair on joint duty:

\[\chem{H_2CO_3} + \chem{H_2O} \rightleftharpoons \chem{HCO_3^-} + \chem{H_3O^+}\]

When outside acid enters blood, \chem{HCO_3^-} catches it to form carbonic acid; when a base enters, carbonic acid donates a proton to neutralize it. Following Chapter 31's Le Chatelier principle, equilibrium shifts back and forth to hold blood pH firmly between 7.35 and 7.45. Your cells function normally only within this narrow 0.1-unit window: enzyme shapes and ion charges are sensitive to acidity. Blood pH below 7.2 or above 7.6 can threaten life. Carbonic acid also has an exit: decomposing into \chem{CO_2}, which the lungs exhale. Breathing depth therefore helps regulate acidity; holding your breath for a minute makes blood slightly more acidic by retaining \chem{CO_2}. Respiration and buffering work as a team; keep that clue for later.

\textbf{Titration} is the chemical detective's quantitative craft. How much acetic acid is in a bottle of vinegar? Add caustic soda of known concentration dropwise, monitoring with an indicator or pH meter. Phenolphthalein is colorless in acid and turns pink above about pH 8.2. The first hint of pink indicates that the acid has been neutralized with not one drop too many or too few. Record the base volume used and calculate from the reaction's 1:1 ratio to find the total acetic acid. Titration's essence is \textbf{inferring an unknown amount from a precise volume and a known concentration}. Chapter 35 extends this accounting by reaction ratios into a broader chemical detective toolkit.

\section{Three Open Accounts: Teeth, Oceans, Blood}

Take your pH perspective outside the classroom, and three accounts are unfolding around you.

The first is in your mouth. Tooth enamel consists mainly of hydroxyapatite, a mineral containing calcium, phosphate, and hydroxide. Acid dissolves it: below an oral pH of about 5.5, enamel starts losing minerals. Mouth bacteria ferment sugar into acid, so sugar itself is not the real culprit but the acid it feeds. Cola is worse, bringing its own pH of 2.5 plus sugar. Fortunately, saliva is a buffer that restores pH toward safety, and fluoride helps enamel remineralize. Frequently sipping carbonated drinks bathes teeth in acid all day. Damage depends on \textbf{the total exposure dose}, again returning to the dose principle.

The second account spans the ocean. Atmospheric \chem{CO_2} continually dissolves in seawater, following the dry-ice route from the opening: form carbonic acid, then donate protons. Since the Industrial Revolution, fossil-fuel burning has raised atmospheric \chem{CO_2} by about half, while surface seawater pH has fallen from about 8.2 to 8.1, the approximately $26\%$ hydrogen-ion increase calculated earlier. The problem is that corals and shellfish build with calcium carbonate, \chem{CaCO_3}, while the extra \chem{H^+} captures carbonate, \chem{CO_3^{2-}}, to form \chem{HCO_3^-}, taking bricks away from the building site. This is ``ocean acidification'': the ocean does not actually become ``acidic,'' since pH remains above 8, but movement toward acidity raises costs for life dependent on calcium carbonate.

The third account is in your blood vessels and never closes. Every step you run makes muscles produce lactic acid and \chem{CO_2}; blood buffers, lungs, and kidneys work rotating shifts to hold pH near 7.4. They adjust Chapter 31's chemical equilibrium and defend Chapter 32's acid--base balance. Volume II's metabolism will show how this is written into life's operating system.

\section{Back to Lemon and Soap}

Now to fulfill the opening promise.

Lemon's sourness results from citric acid stored in fruit cells donating protons in water, and \chem{H_3O^+} meeting receptors on your tongue. Your tongue is not ``tasting lemon'' but measuring pH. This sensor is so ancient that all vertebrates use it, because body-fluid acid--base balance is a life-or-death reading for all life. This also solves a small puzzle: if citric acid is weak, why is lemon so sharply sour? Because it is \textbf{concentrated}: abundance brings lemon juice down to pH 2.4, again showing that strength and concentration are separate.

Soap's slipperiness is a double trick. Soap is itself a fatty-acid salt made by saponifying fat with base. Its molecules have water-loving and oil-loving ends, meeting Chapter 22's intermolecular forces again, and disperse grease. But some of the slippery skin sensation comes from residual alkalinity actually saponifying your surface oils on the spot. A small part of the soap you feel comes from your own oils. This explains why strongly alkaline cleaners leave hands dry and tight: skin oils have been stripped away. Slipperiness is not soap's personality, but the sound of base working on your skin.

\section{Protons Change Hands; Electrons Do Too}

Looking back, this chapter discussed just one action: a positively charged particle, the proton, passes from one molecule to another. Acids donate, bases accept, neutralization returns things to their place, and buffers are relay teams always on standby.

But protons are not the only particles that change hands in chemistry. Another, more precious and dangerous particle constantly moves between molecules: the electron. Iron rusts by transferring electrons to oxygen; batteries discharge as electrons are forced to detour through an external circuit. Breathing, eating, and burning a lump of sugar all ultimately involve electrons flowing from one substance to another. Proton transfer defined acids and bases; electron transfer defines the next chapter's protagonist: oxidation and reduction.

\begin{tiaozhan}
Skipping this does not break the main thread. Strictly, the brackets in pH's definition denote not concentration but \textbf{activity}, a particle's ``effective concentration'' in a crowded solution. Activity approximately equals concentration in dilute solutions, so using concentration directly at school level causes no trouble. But in concentrated solutions or an ``ion soup'' such as seawater, ions constrain one another and activity becomes appreciably lower than concentration. This is one reason precise seawater pH measurement requires special electrodes and calibration curves.

Derive neutrality's temperature drift once more. Define water's ion product as $K_w=[\chem{H^+}][\chem{OH^-}]$. Experiments give $K_w\approx1.1\times10^{-15}$ at $0\,^{\circ}\mathrm{C}$, $K_w=1.0\times10^{-14}$ at $25\,^{\circ}\mathrm{C}$, and $K_w\approx5.1\times10^{-13}$ at $100\,^{\circ}\mathrm{C}$. In pure water, the two concentrations are equal, so $[\chem{H^+}]=\sqrt{K_w}$. At $100\,^{\circ}\mathrm{C}$ this gives $[\chem{H^+}]\approx7.1\times10^{-7}$ mol/L and pH $=7-\lg 7.1\approx6.14$; at $0\,^{\circ}\mathrm{C}$, pH $\approx7.48$. Three temperatures, three ``neutral pH'' values, yet water is strictly neutral at all three. $K_w$ rises with temperature because ionization is endothermic and equilibrium moves forward. This is not a memorized table but a quantitative victory for Le Chatelier's principle.
\end{tiaozhan}

\section*{Try It Yourself}

\begin{sikaoti}
\item Draw particle diagrams of two beakers containing concentrated hydrochloric and concentrated acetic acid, assuming equal total numbers of solute molecules. Use circles for unionized molecules and paired charged balls for ions to show the different particle types. Note that particle sizes and numerical proportions are schematic, not real ratios.
\item Someone separately dissolves \ud{0.1}{mol} of hydrochloric acid and \ud{0.1}{mol} of acetic acid in \ud{1}{L} of water. Which solution is more concentrated? Which is more acidic? Which releases more \chem{CO_2} after reacting fully with excess baking soda? Hint for the last question: use Le Chatelier's principle. As weak acid is continually consumed, how does equilibrium shift?
\item ``Normal'' rainwater has pH about 5.6, not 7. Using this chapter and the preceding section, explain why rain is weakly acidic even without pollution. How many times higher is hydrogen-ion concentration in pH 5.6 rain than in pure water at $25\,^{\circ}\mathrm{C}$?
\item Pure water is heated from $25\,^{\circ}\mathrm{C}$ to $80\,^{\circ}\mathrm{C}$. Does its pH rise, fall, or stay constant? Has it become acidic? Explain using the relationship between $[\chem{H^+}]$ and $[\chem{OH^-}]$, and identify the error in saying ``pH 7 means neutral.''
\item Draw a material-flow diagram for someone who loses much stomach hydrochloric acid through severe vomiting. Using arrows and conjugate acid--base pairs, predict which way blood pH shifts and which way its carbonic acid buffer moves to ``fight the fire.'' Hint: this real clinical problem is metabolic alkalosis. Include at least \chem{H_2CO_3}, \chem{HCO_3^-}, and \chem{CO_2} in your diagram.
\end{sikaoti}
